\documentclass[11pt]{article}

\usepackage[utf8]{inputenc}
\usepackage[T1]{fontenc}
\usepackage{amsmath,amssymb,amsfonts,amsthm}
\usepackage{booktabs}
\usepackage{hyperref}
\usepackage{geometry}
\usepackage{enumitem}
\geometry{margin=1in}

\title{\textbf{Feature--Entity Joint Localization}\\
\large via Conditional Mean Differences\\
\large (Non-causal; Not Subgroup Analysis)}
\author{TencentGR / multimodal attribution prototype}
\date{2026-09-19}

\newtheorem{definition}{Definition}
\newtheorem{proposition}{Proposition}
\newtheorem{remark}{Remark}

\newcommand{\E}{\mathbb{E}}
\newcommand{\R}{\mathbb{R}}
\newcommand{\ind}{\mathbf{1}}
\newcommand{\norm}[1]{\left\|#1\right\|_2}

\begin{document}
\maketitle

\begin{abstract}
We localize period-to-period distribution shift by shrinking the
\emph{support} of edges with an entity key, using only
\textbf{conditional mean} contrasts between windows \(W_1\) and \(W_2\).
The primitive object is an entity--feature matrix of mean differences
\(D\in\R^{U\times d}\). Columns of \(D\) (and the parent mean gap
\(\delta\)) guide which features matter; rows of \(D\) decide whether
to drill to a finer entity support. Feature--entity joint rules are
declared as ordered compositions \(\mathrm{F}\!\to\!\mathrm{E}\) or
\(\mathrm{E}\!\to\!\mathrm{F}\) on the same \(D\). The protocol is
modality-agnostic and \emph{non-causal}: \(W\) is a period indicator,
not a treatment. It is \emph{not} subgroup analysis: we retrieve a
stable localization kernel, we do not estimate heterogeneous effects.
\end{abstract}

\tableofcontents
\vspace{1em}

%----------------------------------------------------------------------
\section{Claim and justification (read first)}
\label{sec:justify}

\paragraph{What we optimize.}
Given edges observed in two periods, find a smaller support
\(S_{t+1}\subset S_t\) on which the period mean-shift is concentrated,
at a favorable capture-to-shrinkage rate (Sec.~\ref{sec:rate}).

\paragraph{What we do \emph{not} claim.}
No ATE/CATE; no ``these merchants have larger treatment effects'';
no confirmatory hypothesis test on drill gates. Window label \(W\) is
\emph{period}, not treatment.

\paragraph{Why this looks like subgroup analysis---and why it is not.}
Procedurally we shrink support in a data-dependent way, which is
\emph{isomorphic in control flow} to post-hoc subgrouping. The
\emph{target functional} differs:

\begin{center}
\begin{tabular}{@{}lll@{}}
\toprule
 & Subgroup analysis & This protocol \\
\midrule
Object & \(\tau(S)=\E[Y(1)-Y(0)\mid S]\) &
  \(\delta(S)=\mu_{S}^{(1)}-\mu_{S}^{(0)}\) on features \\
Role of \(S\) & Inferential subgroup & Retrieval / delivery kernel \\
``Better'' & Larger / significant effect &
  Sharper row spectrum / lower residual \(R\) \\
Justification & Identification + selective inference &
  Stable retrieval + business cap + descriptive gates \\
\bottomrule
\end{tabular}
\end{center}

\paragraph{Required write-time justification (J1--J8).}
Reports and code docs must state: (J1)~task $=$ period shift
localization; (J2)~entity key atoms; (J3)~column vs row split;
(J4)~stop $=$ flat rows / business cap / rejected \(R\), not
``nonsignificant effect''; (J5)~deliver \(\pi\)-stable kernels;
(J6)~any Lift is enrichment \(\bar s_S/\bar s_C\); (J7)~optional
holdout evaluation; (J8)~business granularity cap. Omitting these
causes the reader to parse the pipeline as subgroup fishing.

\begin{remark}[Conditional mean is enough for the gate]
Drill / stop uses only conditional means (and aggregates of an optional
once-fit edge score). Full MMD leave-one-entity recomputation is
unnecessary for the decision; MMD may audit top candidates.
\end{remark}

%----------------------------------------------------------------------
\section{Setup: periods, edges, entity keys}
\label{sec:setup}

Let each row \(i=1,\ldots,n\) be an edge (or interaction) with
standardized features \(X_i\in\R^d\) (scaler fit on \(W_1\) only),
period label \(W_i\in\{0,1\}\) (\(0=W_1\), \(1=W_2\)), and one or more
entity keys. A \emph{key} \(k\) is a column (e.g.\ merchant, user,
category); a \emph{value} \(v\) is a level of that key. Write
\(e_i^{(k)}\) for the value of key \(k\) on edge \(i\).

Fix a parent support \(S\subseteq\{1,\ldots,n\}\) (all edges, or a
previous kernel). Mass gate for key \(k\):
\begin{equation}
  U_{\mathrm{ok}}^{(k)}
  =\Bigl\{
    v:\ 
    \#\{i\in S: e_i^{(k)}=v,\,W_i=0\}\ge n_{\min},\;
    \#\{i\in S: e_i^{(k)}=v,\,W_i=1\}\ge n_{\min}
  \Bigr\}.
\end{equation}

%----------------------------------------------------------------------
\section{Primitive: conditional means and the joint matrix \(D\)}
\label{sec:cmean}

\begin{definition}[Window conditional means]
For support \(S\) and entity value \(v\) under key \(k\),
\begin{align}
  \mu_S^{(w)}
  &=\E\bigl[X \mid i\in S,\, W=w\bigr]
  \in\R^d, \\
  \mu_{v}^{(w)}
  &=\E\bigl[X \mid e^{(k)}=v,\, i\in S,\, W=w\bigr]
  \in\R^d.
\end{align}
Empirical means replace \(\E\) in code (one group-by pass).
\end{definition}

\begin{definition}[Parent gap and entity--feature joint matrix]
\begin{align}
  \delta(S)
  &=\mu_S^{(1)}-\mu_S^{(0)}
  \in\R^d,
  \label{eq:delta} \\
  D^{(k)}_{v,:}
  &=\mu_v^{(1)}-\mu_v^{(0)}
  \in\R^d,
  \qquad v\in U_{\mathrm{ok}}^{(k)}.
  \label{eq:D}
\end{align}
Stacking rows yields \(D^{(k)}\in\R^{|U_{\mathrm{ok}}|\times d}\):
\textbf{rows $=$ entity values, columns $=$ features}. All feature and
entity decisions in this note are projections of the same \(D^{(k)}\)
and \(\delta(S)\).
\end{definition}

\paragraph{Why conditional mean.}
It is \(O(nd)\) after standardization, aligns with a mean embedding
view of shift, and makes leave-one-\emph{entity} updates additive in
the embedding (no \(|U|\) full MMD refits). Higher-order discrepancies
(full MMD) remain optional audits.

%----------------------------------------------------------------------
\section{Column view: feature guidance}
\label{sec:cols}

\begin{definition}[Feature energy shares]
\[
  \mathrm{share}_j(S)
  =\frac{\delta_j(S)^2}{\norm{\delta(S)}^2+\varepsilon}.
\]
Let \(J^\star(S)=\mathrm{TopShare}\bigl(\mathrm{share}(S)\bigr)\).
\end{definition}

\begin{definition}[Feature gates (not drill)]
\begin{description}[leftmargin=2.2em]
  \item[G0] \(\norm{\delta(S)}<\varepsilon_\delta\):
    negligible mean shift; skip FSDS at this support.
  \item[G1] large \(\norm{\delta}\) and concentrated share:
    if we \emph{stop} here, run FSDS once on \(S\) and report \(J^\star\).
  \item[G2] large \(\norm{\delta}\) but flat share:
    weak feature tip; FSDS optional.
\end{description}
\end{definition}

\noindent
Allowed conclusions: ``shifted / not'', ``watch features \(J^\star\)'',
``FSDS if stop''. Forbidden: inferring drill from G0--G2 alone.

%----------------------------------------------------------------------
\section{Row view: drill / stop on an entity key}
\label{sec:rows}

\begin{definition}[Row spectrum]
\[
  r_v^{(k)}
  =\norm{D^{(k)}_{v,:}},
  \qquad
  a_v^{(k)}
  =\frac{\langle D^{(k)}_{v,:},\,\delta(S)\rangle}
        {\norm{D^{(k)}_{v,:}}\,\norm{\delta(S)}+\varepsilon}.
\]
Tail and CV of \(\{r_v\}\) (over \(U_{\mathrm{ok}}\)) measure whether
shift mass is entity-heterogeneous under key \(k\).
\end{definition}

\begin{definition}[Drill gates]
\begin{description}[leftmargin=2.8em]
  \item[D-stop] low Tail/CV (and \(a_v\approx 1\) for most \(v\)):
    homogeneous entities \(\Rightarrow\) do not refine key \(k\).
  \item[D-drill] high Tail or CV:
    \(K=\mathrm{Top}(r^{(k)})\cap\{\pi_v\ge\pi^\star\}\)
    (stability \(\pi\) from probes/bootstrap).
  \item[D-check]
    residual ratio on means
    \[
      R(K)
      =\frac{\norm{\delta(S\setminus K)}^2}{\norm{\delta(S)}^2+\varepsilon}.
    \]
    Accept \(K\) iff \(R(K)\le R^\star\); else reject false tip.
\end{description}
\end{definition}

\begin{definition}[Capture and shrinkage rates]
\[
  \alpha(K)=\frac{|S_K|}{|S|},
  \qquad
  \gamma(K)=1-R(K),
  \qquad
  \mathrm{eff}(K)=\frac{\gamma(K)}{\alpha(K)+\varepsilon}.
\]
Localization prefers high \(\gamma\) at moderate \(\alpha\), with stable
\(\pi\). Early stop iff further refinement yields \(\gamma\approx 0\).
\end{definition}

%----------------------------------------------------------------------
\section{Feature--entity joint prototypes}
\label{sec:joint}

The joint logic is an \emph{ordered} use of columns then rows (or the
reverse) on the same conditional-mean matrix. Order must be declared;
silent unions of two orders are forbidden.

\subsection{Path \(\mathrm{F}\to\mathrm{E}\) (default)}
\label{sec:FE}

\begin{enumerate}[leftmargin=1.6em]
  \item Compute \(\delta(S)\) and \(J^\star(S)\) (Sec.~\ref{sec:cols}).
  \item Restrict columns:
    \(\widetilde D_{v,:}=D^{(k)}_{v,\,J^\star}\).
  \item Set \(r_v=\norm{\widetilde D_{v,:}}\) and apply D-stop / D-drill /
    D-check (Sec.~\ref{sec:rows}).
\end{enumerate}
\paragraph{Role.} Denoise the row spectrum when irrelevant coordinates
dilute \(\norm{D_{v,:}}\). \(J^\star\) comes from parent \(\delta\), not
from a fresh FSDS fit.

\subsection{Path \(\mathrm{E}\to\mathrm{F}\)}
\label{sec:EF}

\begin{enumerate}[leftmargin=1.6em]
  \item Full-dimension \(r_v=\norm{D^{(k)}_{v,:}}\) \(\to\) kernel \(K\).
  \item Recompute \(\delta(K)\) and \(J^\star(K)\) on edges in \(K\).
  \item Explain / FSDS only if stopping at \(K\).
\end{enumerate}
\paragraph{Role.} When parent \(\delta(S)\) is washed out but a few
entities carry large \(D_{v,:}\).

\subsection{Cell heatmap (explanation only by default)}
\[
  C_{vj}=\bigl|D^{(k)}_{vj}\bigr|.
\]
Large cells suggest ``value \(v\) drifts on feature \(j\)''. Do not use
\(U\times d\) cells as the default drill gate; if used, restrict to
\(J^\star\times\mathrm{Top}(r)\) and label exploratory.

\begin{proposition}[Same object, two projections]
Feature guidance and entity drill are not two estimators: they are the
column energy of \(\delta\) (or of \(D\)) and the row energy of \(D\)
(optionally column-restricted). Joint paths only choose
\emph{which projection is applied first}.
\end{proposition}

%----------------------------------------------------------------------
\section{Multiple entity keys}
\label{sec:keys}

Nested keys (tree): \(\mathrm{merchant}\to\mathrm{user}\to\mathrm{order}\).
Sibling keys (same \(S\)): e.g.\ category vs geo vs merchant. For each
key \(k\) compute \((\alpha^{(k)},\gamma^{(k)},\pi^{(k)})\) and select
by efficiency and stability; report \(\texttt{drill\_key}\) explicitly.
Intersecting kernels across keys is optional and must be labeled.

%----------------------------------------------------------------------
\section{Optional side score (same partition)}
\label{sec:po}

Let \(s\in\R^n\) be an edge-level period score from a \emph{single} fit
(PO-risk style). Align to the same key:
\[
  g_v=\mathrm{mean}(s\mid e^{(k)}=v),\qquad
  \eta^2=\frac{\mathrm{Var}_v(g_v)}{\mathrm{Var}_i(s)}.
\]
Concordant sharpness with \(r_v\) increases confidence; discordance is
reported, and by default does not overturn D-stop. Forbidden:
\(\mathrm{Var}_x(\hat\tau(x))\) as subset heterogeneity; refitting PO
per entity.

%----------------------------------------------------------------------
\section{Decision table}
\label{sec:table}

\begin{center}
\begin{tabular}{@{}cccc@{}}
\toprule
\(\norm{\delta}\) & row spectrum \(r\) & action & FSDS \\
\midrule
small & flat & stop & no \\
small & sharp & rare; check mass & cautious \\
large & flat & stop drill; feature story & once on \(S\) if G1 \\
large & sharp & drill along key to \(K\) & once on final \(K^\star\) \\
\bottomrule
\end{tabular}
\end{center}

\noindent
Business cap truncates delivery depth independently of \(\gamma\).

%----------------------------------------------------------------------
\section{What else sits on the same backbone}
\label{sec:else}

Default: nested path + early stop, sibling-key comparison,
\(\mathrm{F}\to\mathrm{E}\), optional \(\mathrm{E}\to\mathrm{F}\),
optional holdout \(\gamma\)/overlap, optional peel curve.
Not in the drill product: Causal Forest \(\hat\tau\) variance.
No third methodology---only support shrinkage, multiple keys, and
feature restriction on one conditional-mean matrix \(D\).

%----------------------------------------------------------------------
\section{Minimal algorithm (pseudocode)}
\label{sec:algo}

\begin{verbatim}
# Parent support S, entity key k, features X (W1-scaled), window W
mu_S[w]  = mean(X | S, W=w)
delta    = mu_S[1] - mu_S[0]
share    = delta**2 / (||delta||**2 + eps)
J_star   = top_share(share)                 # feature guidance

for v in values(k) with mass_ok:
    D[v] = mean(X|v,W=1) - mean(X|v,W=0)

# declare joint order:
# F->E:  r[v] = ||D[v, J_star]||
# E->F:  r[v] = ||D[v, :]||  then recompute delta on K

if flat(r): STOP; maybe FSDS(S) if ||delta|| large
else:
    K = top(r) intersect stable_pi
    R = ||delta(S\K)||**2 / ||delta(S)||**2
    if R <= R_star: parent <- edges in K; recurse or deliver
    else: reject tip; STOP at S
# FSDS once on final K_star only
\end{verbatim}

\vspace{1em}
\noindent\textbf{Companion markdown:}
\texttt{docs/tencent\_gr/ATTRIBUTION\_DRILL\_STOP.md}.

\end{document}
